\NumberedItem{4}
\HistoricTitle{\textsc{Gauss} an \textsc{Bessel}. Göttingen, den 18. Dezember 1811.}
\HistoricCitation{[Briefwechsel zwischen \textsc{Gauss} und \textsc{Bessel}. Leipzig 1880, S. 155.]}
\noindent\hspace*{\parindent}Seit einigen Tagen habe ich endlich das \textls[120]{Königsberger Archiv}, das ich mir verschrieben hatte, empfangen. Mit grossem Interesse habe ich Ihren ersten Aufsatz[\textsuperscript{**})], theuerster \textsc{Bessel}, gelesen und den andern[\textsuperscript{***})] wenigstens vorerst durchblättert. Jener und die schönen Hülfsmittel, die er darbietet, das li für grosse Zahlen zu berechnen, sind mir jetzt um so angenehmer, da ich unlängst eine schöne in diesem Jahr zu Deventer erschienene Factorentafel von \textsc{Chernac}[$\dagger$)] bis $1\,020\,000$ erhalten habe, aus welcher ich nach und nach die Primzahlen von Myriade zu Myriade abzählen lassen will, um sie
\par\medskip\noindent\rule{.24\linewidth}{.4pt}\par\smallskip
{\small
\noindent\hspace*{\parindent}[\textsuperscript{**}) \textit{Untersuchung der durch das Integral $\displaystyle\int\frac{dx}{lx}$ ausgedrückten transcendenten Function}, Königsberger Archiv für Naturwissenschaften und Medicin \textbf{1}, 1812, S. 1, \textsc{Fr. W. Bessels} Abhandlungen II, Leipzig 1876, S. 330.]
\par\noindent\hspace*{\parindent}[\textsuperscript{***}) \textit{Über die Theorie der Zahlenfacultäten}, Königsberger Archiv u.s.w. \textbf{1}, 1812, S. 241, \textsc{Fr. W. Bessels} Abhandlungen II, 1876, S. 342.]
\par\noindent\hspace*{\parindent}[$\dagger\dagger$) \textit{Cribrum Arithmeticum sive tabula continens numeros primos etc.}, confecit \textsc{Ladislaus Chernac}, Pannonius, Daventriae 1811, vergl. auch \textsc{Gauss}’ Anzeige dieses Werks, Gött. Gel. Anzeigen 1812, S. 477, Werke II, S. 181.]
\par}
\clearpage\SourceLeaf{373}\ArithmeticHead{366}{L}{ANALYSIS.\quad BRIEFWECHSEL.}
\noindent mit dem Werthe des Integrals
\[
\int\frac{e^x-1}{x}dx
\]
\noindent von $x=0$ an gerechnet zu vergleichen[\textsuperscript{*})]. Sie haben mir den Wunsch zu erkennen gegeben, dass ich diese Aufsätze in unseren Gel[ehrten] Anz[eigen] anzeigen mögte: unsre Freundschaft, lieber \textsc{Bessel}, macht es mir zur Pflicht, ehe ich es thue, mich schriftlich mit Ihnen über einen oder den andern Punct zu unterhalten, wo meine Ansicht nicht ganz mit der Ihrigen übereinstimmt. Nehmen Sie also die folgenden Bemerkungen gütig auf und theilen mir Ihre Meinung darüber eben so freimüthig und offen mit, als ich die meinige. Ich meine aus einer Ihrer gelegentlichen Äusserungen schliessen zu können, dass Ein Grundsatz uns beiden gemein ist: »Es gibt in der Mathematik keine wahren Controversen« und so zweifle ich nicht, dass wir uns durch den wechselseitigen Austausch unserer Ideen schon verständigen werden.

\indent Zuvörderst würde ich jemand, der eine neue Function in die Analyse einführen will, um eine Erklärung bitten, ob er sie schlechterdings bloss auf reelle Grössen (reelle Werthe des Arguments der Function) angewandt wissen will, und die imaginären Werthe des Arguments gleichsam nur als ein Überbein ansieht, oder ob er meinem Grundsatze beitrete, dass man in dem Reiche der Grössen die imaginären $a+b\surd{-1}=a+bi$ als gleiche Rechte mit den reellen geniessend ansehen müsse. Es ist hier nicht von praktischem Nutzen die Rede, sondern die Analyse ist mir eine selbstständige Wissenschaft, die durch Zurücksetzung jener fingirten Grössen ausserordentlich an Schönheit und Ründung verlieren und alle Augenblick Wahrheiten, die sonst allgemein gelten, höchst lästige Beschränkungen beizufügen genöthigt seyn würde. Ich glaube voraussetzen zu müssen, dass Sie im Wesentlichen über diesen Punct mit mir einstimmen, da schon Ihre Erklärungen im Artikel 18.[\textsuperscript{**})] zeigen, dass Sie keineswegs sich den Weg zu Untersuchungen über $\operatorname{li}(a+bi)$ versperren wollen. Was soll man sich nun bei $\int\varphi x{.}\,dx$ für $x=a+bi$ denken? Offenbar, wenn man von klaren Begriffen ausgehen will, muss man annehmen, dass $x$ durch unendlich kleine Incremente (jedes von der Form $a+\beta i$) von
\par\medskip\noindent\rule{.24\linewidth}{.4pt}\par\smallskip
{\small
\noindent\hspace*{\parindent}[\textsuperscript{*}) Vergl. hierzu die Werke II, S. 435 abgedruckte \textit{Tafel der Frequenz der Primzahlen} und den Brief an \textsc{Encke} vom 24. Dezember 1849, ebenda S. 444 ]
\par\noindent\hspace*{\parindent}[\textsuperscript{**}) Siehe \textsc{Fr. W. Bessels} Abhandlungen II, Leipzig 1876, S. 341, 1. Spalte.]
\par}
\clearpage\SourceLeaf{374}\ArithmeticHead{367}{R}{GAUSS AN BESSEL.}
\noindent demjenigen Werthe, für welchen das Integral 0 seyn soll, bis zu $x=a+bi$ übergeht, und dann alle $\varphi x{.}\,dx$ summirt. So ist der Sinn vollkommen festgesetzt. Nun aber kann der Übergang auf unendlich viele Arten geschehen: so wie man sich das ganze Reich aller reellen Grössen durch eine unendliche gerade Linie [dargestellt] denken kann, so kann man das \textls[120]{ganze} Reich \textls[120]{aller} Grössen, reeller und imaginärer Grössen, sich durch eine unendliche Ebne sinnlich machen, worin jeder Punct, durch Abscisse $=a$, Ordinate $=b$ bestimmt, die Grösse $a+bi$ gleichsam repräsentirt. Der stetige Übergang von einem Werthe von $x$ zu einem andern $a+bi$ geschieht demnach durch eine Linie und ist mithin auf unendlich viele Arten möglich.

\indent Ich behaupte nun, dass das Integral $\int\varphi x{.}\,dx$ nach zweien verschiedenen Übergängen immer einerlei Werth erhalte, wenn innerhalb des zwischen beiden die Übergänge repräsentirenden Linien eingeschlossenen Flächenraumes nirgends $\varphi x=\infty$ wird. Diess ist ein schöner Lehrsatz\textsuperscript{*}), dessen eben nicht schweren Beweis ich bei einer schicklichen Gelegenheit geben werde. Er hängt mit schönen andern Wahrheiten die Entwicklungen in Reihen betreffend zusammen. Der Übergang nach jedem Puncte lässt sich immer ausführen, ohne jemals eine solche Stelle, wo $\varphi x=\infty$ wird, zu berühren. Ich verlange daher, dass man solchen Puncten ausweichen soll, wo offenbar der ursprüngliche Grundbegriff von $\int\varphi x{.}\,dx$ seine Klarheit verliert und leicht auf Widersprüche führt. Übrigens ist zugleich hieraus klar, wie eine durch $\int\varphi x{.}\,dx$ erzeugte Function für einerlei Werthe von $x$ mehrere Werthe haben kann, indem man nemlich beim Übergange dahin um einen solchen Punct, wo $\varphi x=\infty$[\textsuperscript{**})] entweder gar nicht, oder einmal, oder mehreremale herumgehen kann. Definirt man z. B. $\log x$ durch
\[
\int\frac{1}{x}\cdot dx,
\]
\noindent von $x=1$ anzufangen, so kommt man zu $\log x$ entweder ohne den Punct $x=0$ einzuschliessen oder durch ein- oder mehrmaliges Umgehen desselben; jedesmal kommt dann die Constante $+2\pi i$ oder $-2\pi i$ hinzu: so sind die
\par\medskip\noindent\rule{.24\linewidth}{.4pt}\par\smallskip
\noindent\hspace*{\parindent}\textsuperscript{*}) Eigentlich ist hiebei noch angenommen, dass $\varphi x$ selbst eine einförmige Function von $x$ ist, oder wenigstens für deren Werth innerhalb jenes ganzen Flächenraumes nur Ein System von Werthen ohne Unterbrechung der Stetigkeit angenommen wird.
\par{\small
\noindent\hspace*{\parindent}[\textsuperscript{**}) In der Handschrift steht $x=\infty$.]
\par}
\clearpage\SourceLeaf{375}\ArithmeticHead{368}{L}{ANALYSIS.\quad BRIEFWECHSEL.}
\noindent vielfachen Logarithmen von jeder Zahl ganz klar. Kann $\varphi x$ nie für einen endlichen Werth von $x$ unendlich werden, so ist das Integral immer nur eine einförmige Function. Diess ist z. B. der Fall für
\[
\varphi x=\frac{e^x-1}{x},
\]
\noindent so dass
\[
\int\frac{e^x-1}{x}dx
\]
\noindent gewiss eine einförmige Function von $x$ ist, deren Werth durch die \textls[120]{immer} convergirende, also immer einen und nur Einen Sinn habende Reihe dargestellt wird
\[
x+\frac{1}{4}xx+\frac{1}{18}x^3+\frac{1}{96}x^4+\text{etc.}
\]
\noindent Ich wollte, Herr \textsc{Soldner} hätte, da er doch einmal eine neue Function einführen wollte[\textsuperscript{*})], statt seines
\[
\operatorname{li}x=\int\frac{dx}{\log x}
\]
\noindent lieber jene gewählt, da eine einförmige Function immer ohne Vergleich als classischer und einfacher anzusehen ist als eine vielförmige, zumal da $\log x$ selbst schon eine vielförmige Function ist. Es wäre vielleicht gut, auch für
\[
\int\frac{e^x-1}{x}dx
\]
\noindent oder wenigstens für
\[
\int\frac{e^x dx}{x}
\]
\noindent ein eignes Zeichen und Namen einzuführen, um so mehr da bei den Aufgaben aus der Physik [, die] auf $\operatorname{li}x$ führen, gemeiniglich $x$ selbst eine Exponentialgrösse ist. Wenn man die Wahrheiten für \textsc{Soldners} li auf mein
\[
\int\frac{e^x dx}{x}\mathpunct{\raisebox{.35ex}{,}}
\]
\noindent das ich einmal Kürze halber mit $\operatorname{Ei}x$, Exponentiallogarithmen bezeichnen will,
\par\medskip\noindent\rule{.24\linewidth}{.4pt}\par\smallskip
{\small
\noindent\hspace*{\parindent}[\textsuperscript{*}) \textsc{Johann Georg} v. \textsc{Soldner}, \textit{Théorie et tables d’une nouvelle fonction transcendante}, München 1809.]
\par}
\clearpage\SourceLeaf{376}\ArithmeticHead{369}{R}{GAUSS AN BESSEL.}
\noindent übertragen will, so ist die Integration so anzunehmen, dass Ei einer reellen negativen unendlichen Grösse verschwindet.

\indent Ich habe alles dieses vorausgeschickt, um meine Ansicht zu begründen, dass ich \textsc{Euler}n beitreten muss, wenn er sagt[\textsuperscript{*})], dass
\[
\int\frac{dx}{\log x}\mathpunct{\raisebox{.35ex}{,}}
\]
\noindent falls es für $x<1$ reell angenommen werde, für Werthe $>1$ nothwendig imaginäre Werthe erhalte. Ihre Differenz von den reellen, welche \textsc{Mascheroni}[\textsuperscript{**})], \textsc{Soldner} und Sie ihm beilegen, ist $=\pi i$ oder $3\pi i$ oder $5\pi i$ etc. Den Übergang durch $x=1$ statuire ich nicht; man würde auf ganz ähnliche Art beweisen können, dass $\log-x=\log+x$ (ein Satz, den man gelten lassen kann, wenn man sich auf reelle Grössen einschränkt, der aber sogleich wegfallen muss, wenn man dem Reiche aller Grössen meinem obigen Grundsatze zufolge zwei Dimensionen beilegt.) Machen Sie $\operatorname{li}x$ reell für irgend einen Werth von $x$ zwischen 0 und 1, gut! Aber welche Werthe wollen sie nun
\[
\begin{gathered}
\operatorname{li}(0{,}5+0{,}001i),\quad\operatorname{li}(0{,}6+0{,}001i),\\
\operatorname{li}(0{,}7+0{,}001i),\quad\operatorname{li}(0{,}8+0{,}001i),\\
\operatorname{li}(0{,}9+0{,}001i),\quad\operatorname{li}(1{,}0+0{,}001i),\\
\operatorname{li}(1{,}1+0{,}001i)\ \text{etc. bis}\ \operatorname{li}(1{,}5+0{,}001i)
\end{gathered}
\]
\noindent beilegen? Imaginäre ohne Zweifel, aber sie sollen doch dem Gesetze der Stetigkeit folgen, nirgends ein Sprung ex abrupto seyn? Gehen Sie dann von $\operatorname{li}(1{,}5+0{,}001i)$, indem Sie den imaginären Theil $0{,}001i$ auf 0 abnehmen lassen, [zu $\operatorname{li}1{,}5$] über, so kommen sie \textls[120]{gewiss nicht} auf einen reellen Werth von $\operatorname{li}1{,}5$, sondern auf einen, welchem $-\pi i$ anhängt\quad Bei dem, was Sie zum Beweise gegen \textsc{Euler} und mich beibringen, finde ich zu erinnern 1) dass Sie sagen, soll $\operatorname{li}x$ im ganzen Umfange reell werden, so muss man etc., allein das Sollen hängt ja nicht von uns ab, wenn die Stetigkeit nicht ohne Ursache aufgehoben werden darf, und es soll ja eben bewiesen werden, dass es im ganzen Umfange reell genommen werden darf. 2) Allerdings ist
\[
\int\frac{dy}{y}
\]
\par\medskip\noindent\rule{.24\linewidth}{.4pt}\par\smallskip
{\small
\noindent\hspace*{\parindent}[\textsuperscript{*}) \textit{Institutiones calculi integralis} I, 1768, § 228, L. \textsc{Euleri} Opera omnia, ser. I, vol. 11, S. 128.]
\par\noindent\hspace*{\parindent}[\textsuperscript{**}) \textsc{Lorenzo Mascheroni}, \textit{Adnotationes ad calculum integralem Euleri}, Ticini, 1790, Adnotatio I, wieder abgedruckt in L. \textsc{Euleri} Opera omnia, ser. I, vol. 11, siehe dort S. 432.]
\par}
\noindent X1.\hfill47\par
\clearpage\SourceLeaf{377}\ArithmeticHead{370}{L}{ANALYSIS.\quad BRIEFWECHSEL.}
\noindent sowohl $\log y$ als $\log-y$, aber nie beides zugleich, sondern das erstere, wenn man das Integral von $y=1$ anfangen lässt, das zweite, wenn es von $y=-1$ anhebt; das zweite Integral ist eben so wie das erste in dem allgemeinen $\log y+C$ begriffen, wenn man das einemal $C=0$, das zweite mal $C=\pm\pi i$ oder $\pm3\pi i$ etc. setzt. Sehr wahr aber ists, dass \textsc{Eulers} Bemerkung insofern einer Berichtigung bedarf, als keineswegs $C$ unendlich wird, wenn das Integral von $z=0$ anheben soll[\textsuperscript{*})]. — Meiner Ansicht nach darf man also nicht setzen
\[
\int\frac{dx}{\log x}=C+l(\pm lx)+\text{etc.},
\]
\noindent sondern muss sich entweder zu $l(+lx)$ oder zu $l(-lx)$ erklären, aber nur zu einem bestimmt entschliessen. —

\indent Übrigens glaube ich, dass die Ausdehnung der Untersuchungen auf imaginäre Argumente zu höchst interessanten Resultaten Anlass geben wird. Doch mögte ich aus den oben angeführten Gründen lieber die Function
\[
\int\frac{e^x-1}{x}dx
\]
\noindent als
\[
\int\frac{dx}{\log x}
\]
\noindent wählen, weil ich vermuthe, dass erstere concinnere Resultate geben wird. So zum Beispiel, möchte ich sehr gern wissen, ob jene Function oder, was dasselbe ist, die Reihe
\[
x+\frac{1}{4}xx+\frac{1}{18}x^3+\text{etc.}
\]
\noindent für gewisse endliche Werthe von $x$ von der Form $a+bi$ wol 0 werden kann. Mit Gewissheit kann ich es noch nicht behaupten, ob wol es mir sehr wahrscheinlich ist. Gibt es solche Werthe (dann gewiss unendlich viele), so werden diess sehr merkwürdige Grössen seyn und die ganze Reihe wird sich in unendliche Factoren der Form $(1+2ax+\beta xx)$ zerlegen lassen[\textsuperscript{**})]. —

\indent Von ein Paar andern Kleinigkeiten kann ich diesmal nur einige Worte hinzusetzen. Ich hätte gewünscht, dass Sie etwas umständlicher (pag. 6[\textsuperscript{***})] am
\par\medskip\noindent\rule{.24\linewidth}{.4pt}\par\smallskip
{\small
\noindent\hspace*{\parindent}[\textsuperscript{*}) Siehe \textit{Institutiones calculi integralis} I, a. a. O.]
\par\noindent\hspace*{\parindent}[\textsuperscript{**}) Vergl. den S. 374 folgenden Brief an \textsc{Bessel} vom 5. Mai 1812.]
\par\noindent\hspace*{\parindent}[\textsuperscript{***}) \textsc{Fr. W. Bessel}, Abhandlungen II, 1876, S. 331, 1. Spalte.]
\par}
\clearpage\SourceLeaf{378}\ArithmeticHead{371}{R}{GAUSS AN LAPLACE.}
\noindent Ende der Note) gezeigt hätten, wie die \textls[120]{allgemeine Gültigkeit} des Gesetzes des continuirlichen Bruches \textls[120]{nach Ihrem angedeuteten Verfahren} folge. Die Constante $0{,}577\,{.}\,{.}\,{.}$ hatte ich schon vor längerer Zeit auf 23 Stellen auf doppeltem Wege aus 10 und 20 Gliedern berechnet und von der 20\textsuperscript{sten} an von \textsc{Mascheroni} verschieden gefunden, nemlich $\overset{\scriptscriptstyle18}{0}60653$[\textsuperscript{*})]; ich glaube, dass meine Zahl die richtige ist, habe aber die Rechnungen selbst nicht aufgehoben. — Mit verschiedenen Ihrer Untersuchungen werden Sie in meiner nun bald vollendeten Abhandlung Berührungspuncte finden. — Die Ableitungen Artikel 17.[\textsuperscript{**})], die sich auf divergierende Reihen und divergierende unendliche Producte gründen, also auf einen immer etwas schlüpferigen Weg, möchten doch wol noch einigen Bedenklichkeiten Platz lassen. —\quad{.}\quad{.}\quad{.}\quad{.}\quad{.}\quad{.}\quad{.}
\par\medskip\noindent\rule{.24\linewidth}{.4pt}\par\smallskip
{\small
\noindent\hspace*{\parindent}[\textsuperscript{*}) Siehe \textsc{Mascheroni}, a. a. O. S. 431; die zwanzigste, einundzwanzigste und zweiundzwanzigste Stelle sind bei \textsc{Mascheroni} in der Tat fehlerhaft, dagegen sind von der dreiundzwanzigsten Stelle an die Ziffern richtig, bis zur letzten (zweiunddreißigsten), die um eine Einheit erhöht werden muß. Man wird hiernach bei \textsc{Mascheroni} einen \textls[120]{Druck}fehler vermuten. Die von \textsc{Gauss} hier als 3 angegebene dreiundzwanzigste Stelle lautet richtig 1, vergl. auch \textit{Disquisitiones circa seriem} art. 31, Werke III, S. 154, Text und Fußnote. Die bei \textsc{Mascheroni} fehlerhaften Ziffern hat auch \textsc{Soldner} a. a. O. S. 13 richtig angegeben.]
\par\noindent\hspace*{\parindent}[\textsuperscript{**}) \textsc{Fr. W. Bessel}, Abhandlungen II, 1876, S. 340, 1. Spalte.]
\par}
